\section{Discussion}
\label{sec:discussion}

This study separated two questions that the end-of-stay benchmark task conflates: how well a deterioration model discriminates, and when its score could be used.
Three findings follow. First, the end-of-stay task cannot support a claim of early warning: it uses rows charted after departure, and quantities that describe how the stay was monitored add to what is
known at arrival. Second, at a stated prediction time two hours after arrival, the recurrent model did not improve on gradient-boosted trees with timing features, and its
gain over a model of the vital-sign sequence alone came mainly from the static profile. Third, conformal sets that meet their marginal coverage rarely contain the
deterioration label for the patients who deteriorate, and the false-alarm burden of class-conditional sets that do contain it is the precision of the score at that sensitivity, which the coverage guarantee does not reveal.

\subsection{What the End-of-Stay Task Measures}

A model that reads the last ten vital-sign rows of a stay is scoring the stay after most of it has happened. ED
duration, the number of measurements, and the per-vital counts, none of which is a physiological value, reached an AUROC of \RVShortAAuroc{} against
\RVEosHybridAuroc{} for the full model, but that figure is not the leaked share of the full model, because much of it restates what is known at arrival: arrival information alone reached \RVIncArr{}. The increment that monitoring adds is smaller, \RVIncWorkOverArr{} (\RVIncWorkOverArrLo{} to \RVIncWorkOverArrHi{}), of which ED duration alone contributes \RVIncDurOverArr{}. Monitoring intensity may record the clinical judgment that the model is meant to anticipate, or may follow deterioration itself, and the design cannot separate the two. What the diagnostic does
establish is that the end-of-stay AUROC combines arrival information, monitoring, and later physiology, and cannot be used to argue for early warning. The audited task comes from an earlier analysis by the author, so this demonstrates a design that the benchmark pipeline makes easy and is not a correction of published studies. The remedy is to tie every input to the time at which it exists, as for diagnostic codes finalized after discharge~\cite{ramadan2025icdleakage}.

\subsection{Discrimination at the Landmark}

Restricting the inputs to what had been recorded two hours after arrival lowered the AUROC of the best model from \RVEosHybridAuroc{} to about \RVAurocXgbT{}. Among the landmark models the differences were small: timing-aware trees
were better than the primary recurrent model by \RVTreesOverHybrid{} (95\% CI \RVTreesOverHybridLo{} to \RVTreesOverHybridHi{}), the recurrent model exceeded logistic regression (\RVCTwoDelta{}), and the largest single gain came from
static features (\RVCThreeDelta{}). This agrees with the MIMIC-IV-ED benchmark, in which deep sequence models did not outperform simpler ones~\cite{xie2022mimic4ed}. On the same stays and the same outcome, the two-hour trees exceeded the arrival trees by \RVCommonTreesGain{} (\RVCommonTreesGainLo{} to \RVCommonTreesGainHi{}), so most of the discrimination is
already available at arrival. These statements apply to
this cohort, outcome, and tuning budget (four configurations per model, with the selected recurrent configuration at a corner of the grid); a larger search could narrow or reverse the differences.

The timing features of the tree ensemble are legitimate inputs at $L$, but they describe monitoring practice, and their value may not transfer to a hospital that charts differently. The concern is concrete in these data: in the later years of the training period the first ED vital-sign row came to repeat the triage measurement (\RVRecEqTriC{} of stays in 2011, \RVRecEqTriTest{} in the test period). For logistic regression and the trees it did not materialize: refitting without the copied rows changed the AUROC by \RVDdTreesChange{} for the trees and \RVDdLrChange{} for logistic regression (Supplementary Table~\ref{S-tab:dedup}), and the advantage of the trees over logistic regression stayed near \RVOrigTreesOverLr{} in every variant, so the timing features do not explain it. The recurrent models were not refitted, but the stored Hybrid GRU-D models changed by \RVGrudHybridChange{} in AUROC when the copied rows were removed from the test stays, well below the \RVCOneDelta{} of the pre-specified contrast with the trees; scoring without a refit does not show what models trained without the copied rows would have learned. The finding bears on any landmark or chronologically split model built on MIMIC-IV-ED that uses the number, timing, or mask of ED vital-sign rows, because a measurement ``by $L$'' means something different before and after the change; such studies should report the share of copied rows by period and a sensitivity analysis without them, which here left the classical models unchanged but need not do so for other models or outcomes. NEWS2 performed well below every learned model (Table~\ref{tab:discrimination}), consistent with three handicaps of this implementation: five of its seven parameters were available, a missing parameter scores as normal, and the score was not designed for a 24-hour intensive-care or death outcome.

\subsection{Coverage Is Not Alarm Precision}

For a rare outcome, the marginal guarantee is met by
sets that almost never contain the deterioration label: class-1 coverage was \RVSetCovOneStaticLac{} for static LAC at a marginal coverage of \RVSetCovStaticLac{}, because the
marginal threshold is set by the many non-events and the model rarely gives an event a risk as high as the threshold requires.
Calibrating within each class closes the gap, as class-conditional validity predicts~\cite{vovk2003mondrian,ding2023classconditional}, and has been used for sepsis mortality~\cite{yang2024sepsisconformal} and reported for imaging~\cite{millar2026classconditionalms}. What a class-conditional set adds is a guarantee on class-1 coverage, not information about precision. At a nominal level of 0.90, the singleton
set $\{1\}$ identified \RVSetSensMondEci{} of the patients with the outcome and had a PPV of \RVSetPpvMondEci{}, so about \RVSetFaMondEci{} in 100 patients received it
without the outcome. The singleton of a static set is issued exactly for the stays whose score exceeds one threshold, so a plain threshold with the sensitivity of static Mondrian-LAC reached a PPV of \RVPrMondSame{} (the set: \RVPrMondPpv{}). The false-alarm burden is therefore the cost of that sensitivity with this model, which the coverage level selects and does not change; precision rose to \RVAFivePpvMondEci{} at a nominal level of 0.95 and \RVAOnePpvMondEci{} at 0.99 only as sets became larger (Supplementary Table~\ref{S-tab:alpha}).

\subsection{Static, Adaptive, and Periodic Calibration}

The static quantile kept marginal coverage close to the target for the primary model (\RVIntCovStatic{}; bootstrap interval \RVStaticBootHybLo{} to \RVStaticBootHybHi{}, which includes the validity threshold of 0.89), but it did not meet the threshold for every model: coverage was \RVOtherXgbStaticCov{} for the trees (\RVStaticBootXgbLo{} to \RVStaticBootXgbHi{}), below the threshold, and \RVOtherLrStaticCov{} for logistic regression. Shifting the intercept of the score on recent labels restored calibration in the large but, with the same shift applied to the calibration scores, left the coverage of the static sets unchanged (\RVIrCovOne{} against \RVIrCovOrig{}). A quantile recomputed on the recalibrated scores gave \RVIrCovTwo{}, like ECI and periodic recalibration; in these data the gain from updating the quantile is small, and its value shows mainly in the drift test, where ECI recovered coverage after the case-mix change and recalibration did not. 

\subsection{Implications for Reporting and Practice}

Three reporting practices follow for studies of ED deterioration models. A study should state the prediction time and restrict both the inputs and the population at risk to it, reporting an end-of-stay AUROC, if at all, next to a workflow-only baseline that costs one logistic regression. A conformal analysis of a rare outcome should report class-specific coverage with the PPV and sensitivity of the singleton set and the false alarms per 100 patients, at more than one coverage level and against the precision--recall curve of the score, because the marginal guarantee constrains none of them; an online calibration method should be compared with periodic recalibration on a stream with a stated order and label delay, as both restored coverage in our data. For practice, class-conditional sets offer a trade-off to be chosen, not an alarm to be trusted: a department would set the coverage level from the false alarms it accepts per patient correctly flagged, which needs clinical input.

\subsection{Limitations}
\label{sec:limitations}

This study has six limitations. First, all results come from one retrospective cohort at one academic medical center, with outcomes ascertained in hospital only, so
deaths after discharge are not captured; neither external nor prospective validation was performed, and both would be needed before any use at the bedside. Second, the
landmark cohort excludes stays that ended within two hours or had an event before then (\RVHtwoTestEligible{} of \RVHtwoTestStays{} test-period stays were eligible), so the findings apply to
patients who are still in the department at the landmark. Third, recording practice changed between the periods: \RVZeroRowsTrainPct{} of eligible training-period stays
have no ED vital sign after triage, against \RVZeroRowsTestPct{} in the test period, and between 2013 and 2016 the first ED vital-sign row came to repeat the triage measurement (Supplementary Table~\ref{S-tab:recording}), so the timing features and the mask and elapsed-time inputs of GRU-D mean different things in the two periods; a change in case mix cannot be excluded as a further contributor; refitting logistic regression and the trees without the copied rows left their discrimination unchanged, the stored recurrent models, scored without the copied rows, changed little, and they were not refitted. Because the split keeps every patient with an earlier stay in the training period, test-period patients also had fewer earlier admissions, and the share with at least one comorbidity fell from \RVCharTrainElixPct{} to \RVCharTestElixPct{}. Calendar time is
recovered only to about one year, so the stream order within a year is arbitrary and rapid drift could be studied only by imposing it, and the triage table has no
timestamp, so the arrival landmark is an operational anchor. Fourth, the modelling budget was limited: four configurations per model, ECI settings (burn-in window and
step-size grid) that differ from the reference experiments of the method, NEWS2 with five of its seven parameters, label delays of at most 1{,}000 arrivals, and no alternative split that places the later stays of returning patients in the test period or ablation of the comorbidity count at the landmark. Fifth,
five seeds capture training variability only, and the bootstrap captures the sampling variability of patients but not of sites. Sixth, several questions were outside the
design: whether the shortcut reflects clinical judgment or monitoring that follows deterioration itself, which threshold probabilities clinicians would accept for the
decision curve, the calibration of the sets within subgroups, clinical benefit, and how clinicians would act on a set that contains both labels. Subgroup calibration of the sets would need more
events per subgroup; clinical benefit and the response to two-label sets would need a cost elicitation with ED clinicians and a prospective evaluation.
